\section{Conclusion and Future Work}
\label{sec:conclusion}

We asked when two common Spark optimizations, shuffle-partition tuning and file compaction,
are worth their one-time cost. In 250 controlled runs, the default of 200 partitions was the
slowest setting on every workload, and a full search for a better value paid for itself after
three runs. Compaction paid for itself at the first run on fragmented input and never on
single-file input. CAO-DF, a simple break-even rule with a noise floor, used these measurements
to choose between doing nothing, tuning and compacting. On four unseen workloads it lowered
estimated total cost by 47\,\% (P1) and 72\,\% (P2) on average, came close to a hindsight
oracle, and avoided the losses of a speed-only strategy for one-off jobs. The main lesson is
that the number of future runs should be part of every optimization decision.

Future work should (1) apply tuning and compaction together, (2) repeat the study with adaptive
query execution enabled, (3) run on a real cluster with cloud storage and larger data, (4) use
billed rather than estimated costs, and (5) add workload features that explain when compaction
saves wall-clock time and not only task time.
